% =============================================================================
% Preamble - tüm paket importları, sayfa düzeni, font, başlık stilleri
% =============================================================================

% --- Encoding & language --------------------------------------------------
\usepackage[T1]{fontenc}
\usepackage[utf8]{inputenc}
\usepackage[english]{babel}

% --- Page geometry --------------------------------------------------------
% A4, 3.5 cm sol cilt payı, 2.5 cm diğer kenarlar
\usepackage[
  a4paper,
  top=2.5cm,
  bottom=2.5cm,
  left=3.5cm,
  right=2.5cm,
  headheight=14.5pt,
  headsep=0.7cm,
  footskip=1.0cm,
]{geometry}

% --- Math packages (must load BEFORE newtxmath to avoid \Bbbk clash) ------
\usepackage{amsmath, amssymb, amsthm}

% --- Fonts (Times New Roman benzeri) --------------------------------------
\usepackage{newtxtext}     % gövde metni
\usepackage{newtxmath}     % matematik fontu (uyumlu, amssymb'den sonra)
\usepackage{microtype}     % blok hizalamada profesyonel görünüm

% --- Line spacing ---------------------------------------------------------
\usepackage{setspace}
\onehalfspacing            % 1.5 satır aralığı (akademik tez standardı)

% --- Paragraph style ------------------------------------------------------
\setlength{\parindent}{1.27cm}   % ilk satır girintisi (Word default)
\setlength{\parskip}{0pt}        % paragraflar arası ekstra boşluk yok

% --- Units & symbols ------------------------------------------------------
\usepackage{siunitx}
\sisetup{
  detect-all,
  group-separator={,},
  group-minimum-digits=4,
}

% --- Graphics & figures ---------------------------------------------------
\usepackage{graphicx}
\graphicspath{{figures/}}
\usepackage{float}         % [H] placement
\usepackage{subcaption}    % alt-figürler
\usepackage{rotating}      % sidewaysfigure (geniş diyagramlar için)

% --- Tables ---------------------------------------------------------------
\usepackage{booktabs}      % toprule/midrule/bottomrule
\usepackage{array}
\usepackage{multirow}
\usepackage{tabularx}
\usepackage{longtable}     % sayfa kıran tablolar

% --- Captions -------------------------------------------------------------
\usepackage[
  font=small,
  labelfont=bf,
  labelsep=period,
  justification=centering,
]{caption}

% --- Colors & code listings -----------------------------------------------
\usepackage{xcolor}
\definecolor{codegray}{rgb}{0.95,0.95,0.95}
\definecolor{codeblue}{rgb}{0.0,0.2,0.6}
\definecolor{codegreen}{rgb}{0.0,0.5,0.0}

% --- Algorithms (pseudo-code) ---------------------------------------------
\usepackage{algorithm}
\usepackage{algpseudocode}

% --- Headers & footers ----------------------------------------------------
\usepackage{fancyhdr}
\pagestyle{fancy}
\fancyhf{}
\fancyhead[L]{\nouppercase{\leftmark}}
\fancyfoot[C]{\thepage}
\renewcommand{\headrulewidth}{0.4pt}
\renewcommand{\footrulewidth}{0pt}

% Chapter ilk sayfasında "plain" stil, sadece sayfa numarası
\fancypagestyle{plain}{%
  \fancyhf{}
  \fancyfoot[C]{\thepage}
  \renewcommand{\headrulewidth}{0pt}
}

% --- Chapter title style (sade, profesyonel) ------------------------------
\usepackage{titlesec}
\titleformat{\chapter}[display]
  {\normalfont\Large\bfseries}
  {\chaptertitlename\ \thechapter}
  {12pt}
  {\LARGE}
\titlespacing*{\chapter}{0pt}{-20pt}{20pt}

% --- TOC / LOF / LOT formatting -------------------------------------------
\usepackage{tocloft}
\renewcommand{\cftchapfont}{\bfseries}
\renewcommand{\cftchappagefont}{\bfseries}
\setlength{\cftbeforechapskip}{6pt}
% LoF/LoT girişleri arası boşluk: çok cümleli caption'ların ikinci sayfaya
% taşmamasını sağlayacak kadar sıkı, fakat sayfa altında geniş boşluk
% bırakmayacak kadar ferah.
\setlength{\cftbeforefigskip}{8pt}
\setlength{\cftbeforetabskip}{8pt}

% --- Bibliography (IEEE style) --------------------------------------------
\usepackage{csquotes}
\usepackage[
  backend=biber,
  style=ieee,
  sorting=none,
  maxbibnames=99,
]{biblatex}
\addbibresource{references.bib}
\setlength{\bibitemsep}{6pt}
\AtBeginBibliography{\normalsize}

% --- Hyperlinks (en sonda olmalı) -----------------------------------------
\usepackage[
  colorlinks=true,
  linkcolor=black,
  citecolor=black,
  urlcolor=blue!60!black,
  pdfborder={0 0 0},
  bookmarksnumbered=true,
]{hyperref}
\usepackage[capitalize,noabbrev]{cleveref}

% --- Custom commands ------------------------------------------------------
\newcommand{\todo}[1]{\textcolor{red}{\textbf{[TODO: #1]}}}
\newcommand{\thesistitle}{On-Device American Sign Language Recognition: A Comparative Study of Image- and Landmark-Based Approaches for Mobile Deployment}
\newcommand{\thesisauthor}{Yusuf Kenan Akg\"un}
\newcommand{\thesisid}{20210702010}
\newcommand{\thesisadvisor}{Bet\"ul Arslan}
\newcommand{\thesisyear}{2026}
\newcommand{\thesisdepartment}{Department of Computer Engineering}
\newcommand{\thesisuniversity}{Yeditepe University}
\newcommand{\thesisfaculty}{Faculty of Engineering}
